%!TEX root = ../template.tex
%% ====================================================================
%% Capitulo 7 -- Discussao e proposta metodologica
%% Primeira versao completa 2026-09-08. Linguagem descritiva, sem rotulos
%% internos de organizacao do trabalho. Proveniencia dos numeros:
%%   estimadores robustos:  outputs/T11/runs/20260830T114345550731Z-363f433da25e
%%   bandas honestas:       outputs/T12/runs/20260830T122854929917Z-459cf49c965c
%%   gate de aceitacao:     outputs/T13/runs/20260830T165208669276Z-b583143cacb9
%%   candidatos robustos:   outputs/T14/runs/20260830T165252925802Z-ad4bf78ea405
%%   fichas por baseline:   outputs/T15/runs/20260830T165318477358Z-3d408fcdae7a
%%   referencia congelada:  outputs/PART_C/C1/runs/20260903T002922246855Z-eba4e32107a0
%%   envelopes fixos:       outputs/PART_C/C6/runs/20260903T003533726541Z-892a18a0e064
%%   injecao de desvios:    outputs/PART_C/C5/runs/20260903T004416308950Z-2c5caf3d5988
%%   verdade conhecida:     outputs/PART_C/C8/runs/20260902T165944891648Z-ae3e8022c6a5
%%                          (stress …T224337596785Z-745ef41d27f7, joint …T181655416848Z-46ca4436f02f)
%% Os lotes da exploracao com referencia congelada que a reinspecao de
%% 2026-09-03 classificou como nao finais (comparacao de esquemas de
%% referencia, janelas de reajuste, injecao sobre estados publicados,
%% agregacao por frequencia, eletricidade) NAO sao citados com numeros;
%% onde entram, entram qualitativamente e assinalados. Depois da
%% republicacao (Task 7 do plano de 2026-09-03), substituir onde marcado
%% com %% PENDENTE-REPUBLICACAO.
%% A matematica de cada instrumento vai para o anexo matematico.
%% ====================================================================

\chapter{Discussão e proposta metodológica}
\label{cha:discussao}

O Capítulo~\ref{cha:diagnostico} estabeleceu que o método de \emph{baselines}
em produção se reproduz exatamente e não é adequado, de forma uniforme, aos
fins que a ISO~50006 lhe atribui. Este capítulo faz três coisas com esse
resultado. Primeiro, lê-o em conjunto: o que os modos de falha têm em
comum, o que a norma exige e não exige, e porque é que a pergunta certa
sobre uma \emph{baseline} de eficiência não é ``prevê bem?'' mas ``revela o
desvio que existe para revelar?''. Segundo, converte o diagnóstico numa
proposta --- um conjunto de linhas orientadoras para construir, aceitar e
manter \emph{baselines} em processo contínuo --- e demonstra-a na própria
unidade: que candidatos da classe linear estática, robustificada e com
bandas que modelam a dependência, passam um critério de aceitação declarado antes de
qualquer resultado; e o que se aprende quando, em vez de reajustar o modelo
todos os anos, se congela uma referência e se mede tudo o resto como desvio
face a ela, incluindo com verdade conhecida por simulação. Terceiro,
devolve a proposta ao instrumento do Capítulo~\ref{cha:plataforma} como
especificação de um bloco de saúde das \emph{baselines}, e declara os
limites de tudo o que aqui se afirma. O capítulo apresenta a proposta e a sua
demonstração, não o método que a suporta: o desenho da avaliação, a classe de
candidatos e o critério de aceitação estão definidos na
Secção~\ref{sec:protocolo-baselines}, e a matemática de cada instrumento
consta do Apêndice~\ref{app:matematica}.

\section{Leitura integrada do diagnóstico}
\label{sec:leitura-integrada}

\paragraph{Reproduzir não é validar.} O primeiro resultado do
Capítulo~\ref{cha:diagnostico} --- a replicação exata das trinta células
vetor--ano --- é a condição de tudo o resto, mas não diz nada sobre a
adequação do modelo: diz que o que se critica é o que está em produção.
Tudo o que se seguiu mostra que um método pode ser reproduzível, documentado
e aplicado com disciplina há seis anos, e ainda assim não servir de forma
uniforme as três funções que lhe são pedidas. Como preditor da relação com a
carga, tem conteúdo demonstrado no fuel gás e nenhum ganho demonstrado nos
vapores; para comparar períodos, \revisto{a igualdade dos declives anuais é rejeitada no vapor de 3~bar sob BH e BY, mas não na sensibilidade com blocos de 60 dias} e, nos outros vetores, de forma que depende do método e de
onde começa o ano; para alarmes, a banda cobre entre $38\%$ e $99{,}7\%$ dos
dias consoante o vetor e o par de anos, e a deteção real não foi medida. Na
energia elétrica o problema é de medição, não do modelo nem da
reestimação: não há resolução diária para avaliar.

\paragraph{Modos de falha distintos, um desenho comum.} Cada vetor falha à
sua maneira --- o fuel gás por deriva paramétrica que a reestimação acomoda;
o vapor de 24~bar por uma reta anual que se comporta como uma média; o de
10~bar por um nível em movimento; o de 3~bar por identificação fraca,
compatível com uma relação que não é linear nem estática --- e a energia
elétrica tem uma limitação de medição, não um modo de falha. A diversidade
importa porque afasta a receita única: um remédio para a deriva do fuel gás
não serve ao 24~bar, onde não há relação com a carga a estabilizar. Mas \revisto{os quatro vetores diários partilham um desenho, sobre o qual se constrói a proposta deste capítulo:} um modelo linear estático, reestimado
anualmente sobre janelas convencionais, aplicado a \revisto{um processo cujo comportamento difere entre anos sem que se tenham identificado as datas, a duração ou as causas dessas diferenças.} O mecanismo é o da Secção~\ref{sec:sintese-modos-falha}:
a reta estimada num ano acomoda nos parâmetros a parte da deriva que se
projeta na carga e na constante, e quando se muda de reta a nova acomoda o
nível que o ano anterior atingiu.

\paragraph{Previsão e deteção são perguntas diferentes.} O diagnóstico mediu
sobretudo previsão --- \emph{skill} face a referências triviais, erro fora
da amostra, cobertura de bandas --- porque é o que se pode medir sem
calendário operacional. Mas a função de uma \emph{baseline} de eficiência
não é prever o consumo de amanhã: é servir de referência estável contra a
qual um desvio se torna visível. As duas exigências entram em tensão, e a
Secção~\ref{sec:poder-preditivo} já a expôs: uma média móvel de trinta dias
prevê o consumo de vapor melhor do que a relação anual com a carga
precisamente porque incorpora o passado recente --- e, ao incorporá-lo,
incorpora também qualquer desvio persistente, tratando-o como o novo
normal. Há, por isso, uma tensão entre adaptar a referência e
preservá-la como termo de comparação. Uma referência fixa pode prever bem
enquanto o regime não muda; um modelo adaptativo pode detetar mudanças,
desde que o desenho as leve para um canal que se monitoriza --- os
parâmetros, o estado, as inovações acumuladas. O que se perde ao adaptar é o
resíduo face a um ponto fixo, e é esse o canal que as políticas avaliadas
neste capítulo usam: o resíduo da referência congelada, sobre o qual corre
um monitor sequencial (Secções~\ref{sec:referencia-fixa}
e~\ref{sec:verdade-conhecida}). O que se propõe não é escolher entre prever
e detetar: é tratar a referência de reporte e a referência de deteção como
dois objetos com funções distintas, quando o método em produção tenta ser
ambos com um só modelo reestimado uma vez por ano.

A Figura~\ref{fig:ref-estado-monitor} separa os três objetos e mostra onde cada um entra.

\begin{figure}[htbp]
  \centering
  \begin{tikzpicture}[x=1cm, y=1cm, font=\scriptsize]
    \begin{scope}[on background layer]
      \fill[corFundo, rounded corners=2pt] (3.55,0.80) rectangle (12.35,2.80);
      \fill[corFG!5, rounded corners=2pt] (3.55,-4.95) rectangle (12.35,-0.15);
    \end{scope}
    \node[ntmini, anchor=north west] at (3.65,2.75)
      {\emph{fixo entre revisões datadas}};
    \node[ntmini, anchor=north west] at (3.65,-0.20)
      {\emph{atualiza-se com cada dia novo}};

    \node[ntcaixa, text width=2.4cm, minimum height=11mm] (dados) at (1.45,1.62)
      {Dados do dia $t$\\[1pt]{\tiny $E_t$, $Q_t$, $Z_t$}};
    \node[ntmini, anchor=north, align=center, text width=3.0cm] at (1.45,0.95)
      {os dados entram uma só vez, pelo resíduo};

    \node[ntdestaque, text width=3.1cm, minimum height=13mm] (ref) at (5.95,1.62)
      {\textbf{Referência} $f_{\mathrm{ref}}$\\[1pt]
       {\tiny fixa, estimada uma vez e datada}};
    \node[ntcaixa, text width=3.2cm, minimum height=13mm] (res) at (10.35,1.62)
      {Resíduo\\[1pt]{\tiny $r_t = E_t - f_{\mathrm{ref}}(Q_t,Z_t)$}};
    \draw[ntsetaf] (dados) -- (ref);
    \draw[ntsetaf] (ref) -- (res);

    \node[ntdestaque, text width=3.1cm, minimum height=13mm] (est) at (5.95,-1.35)
      {\textbf{Modelo temporal}\\[1pt]
       {\tiny atualiza-se a cada dia com o resíduo}};
    \node[ntcaixa, text width=3.2cm, minimum height=13mm] (eio) at (10.35,-1.35)
      {Estado estimado $\hat{x}_t$\\e inovação $\nu_t$};
    \draw[ntsetaf] (res.south) -- ++(0,-0.50) -| (est.north);
    \draw[ntsetaf] (est) -- (eio);

    \node[ntdestaque, text width=3.1cm, minimum height=13mm] (mon) at (5.95,-3.60)
      {\textbf{Monitor}\\[1pt]
       {\tiny estatística e limiar declarados \emph{a priori}}};
    \node[ntcaixa, text width=3.2cm, minimum height=13mm] (inv) at (10.35,-3.60)
      {Sinal $\to$ investigação};
    \draw[ntsetaf] (eio.south) -- ++(0,-0.42) -| (mon.north);
    \draw[ntsetaf] (mon) -- (inv);
    \draw[ntseta] (res.east) -- ++(0.65,0) -- (13.0,-5.35) -- (5.95,-5.35)
      -- (mon.south);
    \node[ntmini, anchor=north west] at (6.15,-5.45)
      {o resíduo também alimenta o monitor, sem passar pelo estado};
  \end{tikzpicture}
  \caption[Referência, estado e monitor como três objetos distintos]%
  {Os três objetos que o método em produção funde num só. A referência
   responde à pergunta ``quanto se esperaria consumir'' e é fixa entre
   revisões datadas; o modelo temporal responde a ``em que estado está a
   unidade agora'' e atualiza-se com cada dia novo; o monitor responde a
   ``isto é anormal'' e obriga a uma estatística e a um limiar declarados
   antes de ver os dados. O alerta não nasce, assim, da escolha do
   normalizador. Complementa a Figura~\ref{fig:desenhos-avaliacao}: aqui
   os componentes, lá a cronologia da avaliação.}
  \label{fig:ref-estado-monitor}
\end{figure}

\paragraph{Três formas de avaliar, três estimandos.} As três avaliações que
este capítulo mobiliza --- a transportabilidade anual do
Capítulo~\ref{cha:diagnostico}, a adjudicação de candidatos da
Secção~\ref{sec:demonstracao} e a referência congelada da
Secção~\ref{sec:referencia-congelada} --- respondem a perguntas diferentes
e estão definidas, com os respetivos estimandos, na
Secção~\ref{sec:protocolo-baselines} e na Tabela~\ref{tab:estimandos}. A
distinção importa para ler o que se segue: o mesmo vetor pode ter leituras
que só parecem contraditórias, como o vapor de 24~bar, ``estável e
calibrado'' na primeira avaliação e, na terceira, a deslocar-se de nível sem
que nenhum normalizador acrescente nada à média de 2020. Não há contradição
--- a primeira pergunta pela relação com a carga ano a ano, a terceira pelo
nível face a um ano fixo.


\section{Implicações para a conformidade normativa}
\label{sec:implicacoes-normativas}

A ISO~50001~\cite{ISO50001_2018} obriga a monitorizar, medir e analisar o
desempenho energético; a ISO~50006~\cite{ISO50006_2014} operacionaliza-o em
indicadores e \emph{baselines}, com variáveis relevantes, fatores estáticos
e normalização; a ISO~50015 e o IPMVP~\cite{ISO50015_2014,IPMVP_2022} tratam
da medição e verificação de poupanças. Como o Capítulo~\ref{cha:enquadramento}
mostrou, \revisto{os referenciais estabelecem objetivos e orientações de validação, mas não especificam um protocolo completo para avaliar a transportabilidade temporal e a deteção sob a dependência observada neste caso} --- e carregam o pressuposto tácito de que a
relação entre consumo e variáveis relevantes é estável no período de
referência e transportável para o de reporte. O Capítulo~\ref{cha:diagnostico}
é, em rigor, um teste desse pressuposto num processo contínuo, e o resultado
é que ele não se sustenta de forma uniforme: \revisto{a igualdade dos declives é rejeitada sob BH e BY no vapor de 3~bar, mas não com blocos de 60 dias; no fuel gás e no de 10~bar, a rejeição também depende do método,}
não é conclusivo no de 24~bar, e na energia elétrica não é testável à
resolução diária.

A consequência normativa é incómoda mas precisa. Um sistema de gestão de
energia pode ter todos os elementos formais --- \emph{baselines} definidas,
variáveis relevantes identificadas, indicadores calculados, revisões
anuais documentadas --- e assentar, em alguns vetores, em modelos sem ganho
demonstrado sobre uma média, cujos parâmetros mudam de sinal entre anos e
cujas bandas de alarme oscilam entre muito poucos e muitos alarmes. As
cláusulas analisadas pedem indicadores e referências adequados e deixam à
organização a demonstração dessa adequação, sem fixar como a fazer; se um
sistema sem essa demonstração passa numa auditoria é uma questão que esta
dissertação não avalia. A ISO~50006 admite a revisão periódica da
\emph{baseline} como método predeterminado, e é essa revisão, aplicada
como rotina anual, que acomoda nos parâmetros o desvio que o indicador
devia mostrar (Secção~\ref{sec:leitura-integrada}).

Daqui decorre a posição desta dissertação: a demonstração de adequação que
a norma deixa à organização precisa de critérios explícitos, e a norma não
os traz. O protocolo de avaliação do Capítulo~\ref{cha:metodos} e as linhas
orientadoras que se seguem são uma proposta concreta desses critérios ---
compatível com o enquadramento normativo, porque não pede nada que a norma
proíba, e uma operacionalização dele, porque transforma em critérios
verificáveis a adequação que a norma pede.

A proposta mede-se também contra os estudos mais próximos do caso
(Secção~\ref{sec:corpus-b-proximos}). \textcite{Beisheim2019} descrevem a
fronteira do modelo e uma regra para o refazer depois de alterações na
instalação, mas não testam fora da amostra. \textcite{Yao2016} propõem
refazer a referência quando o consumo cai três desvios-padrão abaixo dela,
regra que dispara igualmente perante uma melhoria real e perante um modelo
caducado; a linha orientadora de recalibração por evidência separa a
deteção da decisão de recalibrar e certifica o modelo novo só com dados
posteriores. \revisto{\textcite{Hamedi2019} e \textcite{Moghadasi2025} validam fora da amostra: no primeiro, a ordem temporal da partição não está descrita; no segundo, a partição é aleatória,} e
\textcite{Velazquez2013} usam a referência horária de uma unidade de
reformação para alertas sem reportarem o desempenho no conjunto de
validação que separaram. O que esta dissertação acrescenta não é um
componente --- validação temporal, estimação robusta, bandas com dependência
e monitores sequenciais existem na literatura estatística --- mas a
combinação aplicada a uma referência em produção replicada exatamente:
domínio declarado, comparadores triviais, validação encadeada no tempo,
critério de aceitação fixado antes dos resultados e um resultado negativo
publicado como tal.

\section{Linhas orientadoras para a construção e aceitação de \emph{baselines}}
\label{sec:guidelines}

As linhas orientadoras foram formuladas antes da demonstração, revistas em
quatro rondas de revisão adversarial independente
e fixadas na versão que aqui se descreve. Cada uma responde a um modo de
falha do Capítulo~\ref{cha:diagnostico} ou a um defeito apontado na revisão;
nenhuma foi ajustada depois de conhecidos os resultados da
Secção~\ref{sec:demonstracao}. Apresentam-se por ordem de aplicação.

\paragraph{Dados reconciliados como pré-condição.} \revisto{Para aplicação em produção, exige-se a validação de fontes descrita no Capítulo~\ref{cha:plataforma} para o período usado; nesta demonstração, a validação documental cobre apenas janeiro a maio de 2026:} imputações sinalizadas e contadas,
proveniência com \emph{hashes}, e um perfil dos dados --- cobertura,
mecanismo de falta, efeito de qualquer filtro de seleção --- lido e aceite
por uma pessoa antes de qualquer análise. O erro de seleção que invalidou a
análise preliminar desta dissertação (um limiar de carga escrito à mão em vez
de lido da configuração oficial) é o exemplo do que este passo existe para
apanhar.

\paragraph{Domínio definido nas variáveis explicativas, nunca no consumo.} A
\emph{baseline} declara o intervalo de carga em que foi estimada, justificado
por percentis de operação ou limites físicos, e, em separado, um critério de regime
quase estacionário (variação de carga abaixo de um limiar; exclusão das
primeiras horas após um arranque, com duração a fixar pela dinâmica da
unidade --- 48 a 72~horas é um ponto de partida, não um valor validado por
estes dados). Domínio da regressão e regime quase estacionário são condições
diferentes. A seleção faz-se sobre os regressores pré-especificados, nunca
sobre o consumo --- selecionar sobre a variável resposta trunca a
distribuição do erro e enviesa o modelo~\cite{Heckman1979}. A exceção é a
exclusão por erro instrumental comprovado e documentado, que é controlo de
qualidade e não seleção por desempenho; o filtro de consumos nulos herdado
do método de produção é uma exclusão pela resposta, e na unidade estudada
não altera a população (Secção~\ref{sec:modelo-dimensional}). O que a
Secção~\ref{sec:perfil-dados} mostrou --- o limiar oficial separa uma
população operacional distinta, não uma cauda --- torna esta declaração
obrigatória: o domínio de validade existe de qualquer forma; a diferença
está em declará-lo ou descobri-lo. Na aplicação, um dia fora do domínio
recebe o estado ``não avaliado'' e o seu alarme é suprimido: não é uma
anomalia nem um dia normal, é uma pergunta a que o modelo não sabe
responder.

\paragraph{Exclusões por causa, nunca por valor.} Um dia é excluído do treino
por um código de razão pré-registado e ligado ao registo de operações
(paragem, arranque, intervenção, campanha), com a declaração de que usos
essa exclusão afeta --- treino, avaliação de desempenho e avaliação de
alarmes são decisões distintas. Nunca por ser um valor extremo. Fica
declarado o risco residual: um registo operacional pode ele próprio ter sido
motivado por um comportamento energético anómalo, e ``documentado'' não é
sinónimo de ``exógeno''.

\paragraph{Especificação antes de robustificação.} Entre o modelo univariado
na carga e qualquer modelo não linear existe o degrau que a ISO~50006 e o
IPMVP sancionam~\cite{ISO50006_2014,IPMVP_2022}: a regressão linear múltipla com variáveis relevantes
pré-especificadas (densidade da carga, condições de \emph{flash}, sinais de
modo de operação). A Secção~\ref{sec:suficiencia} mostrou o que este degrau
vale na unidade --- consistente num vetor, prejudicial noutros --- e daí a
regra: uma covariável é admitida só por valor preditivo incremental
demonstrado fora da amostra, para a frente no tempo, com um critério que
inclua a banda e não apenas o erro pontual; nunca por catálogo.

\paragraph{Estimação robusta, pré-especificada.} O estimador primário é
declarado antes de ver os dados --- um estimador-M de Huber~\cite{Huber1964,HuberRonchetti2009} com constante
de afinação fixa ($c = 1{,}345$~\cite{HollandWelsch1977}) e escala pela
mediana dos desvios absolutos \ajuste{(Equações~\ref{eq:mat-huber}
e~\ref{eq:mat-mad})} ---
com os mínimos quadrados reportados ao lado, e um limiar de divergência
declarado entre os dois que, se excedido, obriga a investigar: a divergência
pode significar má especificação, não valores atípicos. A alavancagem
trata-se separadamente, porque um estimador de Huber não resiste a pontos de
alavancagem elevada~\cite{RousseeuwLeroy1987}. O estimador de Theil--Sen~\cite{Sen1968}
\ajuste{(Equação~\ref{eq:mat-theil-sen})} serve de referência livre de
distribuição em regressão simples. Um ponto e uma banda robustos são
problemas diferentes, e a linha seguinte trata do segundo.

\paragraph{Bandas condicionais e quantílicas, com três desfechos possíveis.} A banda de alarme
constante $\pm 2\hat\sigma$ assume resíduos independentes e homocedásticos que
a Secção~\ref{sec:diagnosticos-desc} mostrou não existirem. A rota primária
modela explicitamente a dependência residual --- um processo autorregressivo
sobre os resíduos~\cite{BoxJenkins2015}, com desfasamentos civis exatos e
ordem escolhida só no treino \ajuste{(Equação~\ref{eq:mat-centro-ar})} --- e produz intervalos condicionais, a um passo, validados pela
cobertura empírica no ano seguinte. A rota de sensibilidade é a regressão
quantílica~\cite{KoenkerBassett1978} \ajuste{(Equação~\ref{eq:mat-quantil})},
com restrições de não cruzamento~\cite{Chernozhukov2010} e incerteza por \emph{bootstrap}
de blocos --- sobre uma janela plurianual, nunca sobre um ano ($n \approx
300$ dias a $\tau = 0{,}025$ deixa cerca de sete observações em cada cauda).
O ponto essencial é o terceiro desfecho: a avaliação pode concluir que
\emph{nenhuma} rota é suportada para um vetor, e esse veredicto é
reportável e bloqueia o alarme por banda nesse vetor. Uma banda não
validada não é uma banda pior; é a ausência de banda.

\paragraph{Critério de aceitação com validação temporal encadeada.} Um
candidato --- estimador, banda e domínio --- só entra em produção depois de
passar um critério declarado antes de qualquer resultado, avaliado para a
frente no tempo. A formulação original exigia uma partição selada: construção
e afinação em 2020--2023, seleção congelada em 2024, uma única avaliação em
2025 intocado. Foi substituída, por decisão datada, por validação temporal
encadeada --- os mesmos cinco pares anuais do diagnóstico ---, porque 2025 e
o parcial de 2026 já tinham informado o Capítulo~\ref{cha:diagnostico} e
nenhum resultado obtido depois disso pode chamar-se \emph{holdout}. O critério
tem dois membros e aceita um candidato apenas se ambos se cumprem: conteúdo
preditivo --- ganho sobre a média do ano de treino, com o intervalo de
confiança a excluir zero, em pelo menos quatro dos cinco pares
\ajuste{(Equações~\ref{eq:mat-indicadores-gate} e~\ref{eq:mat-gate-historico};
a distinção entre compatibilidade e equivalência está na
Equação~\ref{eq:mat-equivalencia})} (é o que
impede que um modelo vazio, com $R^2 \approx 0$, seja aceite por ter uma
banda razoável, o defeito da primeira versão das linhas orientadoras); e
banda compatível pelo critério adotado --- intervalo da cobertura que
inclui a nominal --- em pelo menos quatro dos cinco pares. O limiar de
quatro em cinco foi fixado antes de qualquer resultado e não se mexe depois;
a Secção~\ref{sec:demonstracao} mostra o que depende dele.

Pré-especificar a regra evita ajustá-la ao resultado, mas não a torna
adequada, e três defeitos ficam declarados. O primeiro é da compatibilidade:
um intervalo largo inclui a nominal com mais facilidade, pelo que a regra
premeia a falta de precisão; o que ela declara é compatibilidade pelo
critério adotado, não calibração demonstrada. O segundo é da contagem: a
regra tolera um par desfavorável em cinco e recusa dois, não resultou de uma
análise dos custos de aceitar ou de rejeitar, e trata como independentes
pares que não o são; os dois membros contam-se separadamente, e um candidato
podia cumprir cada um em quatro pares diferentes e ter os dois em apenas
três --- na demonstração isto não chegou a importar, porque nenhum candidato
chegou a quatro pares nos dois membros. O terceiro é do objeto: exigir ganho
sobre a média exclui por construção uma referência constante, que pode ser
útil para monitorizar um consumo que depende pouco da carga; o critério
aceita normalizadores dependentes da carga, não detetores. Uma versão para
entrada em produção deve fixar com a operação, antes de avaliar qualquer
candidato, uma margem de cobertura, uma precisão mínima e um desfecho
inconclusivo, avaliar equivalência ou não inferioridade~\cite{Lakens2017} a
par da largura da banda e do desempenho de deteção, e simular os erros de
decisão do próprio procedimento.

\paragraph{Recalibração por evidência, com políticas separadas.} Gerar um
alarme, investigar, suspender temporariamente um modelo, recalibrar e
verificar a recalibração são cinco decisões com políticas distintas e
declaradas. A recalibração exige uma mudança detetada --- um sinal de um
monitor sequencial~\cite{Page1954,Lai1995} sobre as inovações, não uma data
no calendário --- mais
uma revisão documentada; mas a revisão pode concluir ``causa desconhecida,
recalibrar com sinalização'', porque exigir causa demonstrada
incondicionalmente atrasaria a adaptação exatamente quando o sistema muda.
Os mesmos resíduos nunca servem para declarar a mudança e para certificar o
modelo recalibrado: a certificação usa só dados posteriores. É esta linha
que substitui a recalibração anual de rotina, e é a resposta direta ao
mecanismo da Secção~\ref{sec:leitura-integrada}.

\paragraph{Mínimos de dados e governança.} Não se estima uma \emph{baseline}
se o número de observações está abaixo de um mínimo declarado (formalizando
o parâmetro que o método de produção já tem), se o treino cobre menos do que
uma fração declarada do domínio de aplicação, ou se a fração imputada excede
um teto. A governança --- versões, registo de alterações, aprovação nominal,
revisão agendada --- e uma ficha de uma página por \emph{baseline} (domínio,
janela, exclusões com códigos de razão, estimador, método de banda e o seu
desfecho, resultados da aceitação, \emph{hash} de proveniência) fecham o
conjunto. A ISO~50001 é uma norma de sistema de gestão; sem governança, as
linhas orientadoras não são adotáveis.

\section{Demonstração na unidade}
\label{sec:demonstracao}

As linhas orientadoras foram demonstradas nos quatro vetores diários da
unidade, sobre os mesmos cinco pares anuais 2020${\to}$2021 a
2024${\to}$2025 do Capítulo~\ref{cha:diagnostico}, com constantes fixadas a
priori ou afinadas apenas dentro do ano de treino, conclusões só por
exclusão do intervalo de confiança por \emph{bootstrap} de blocos móveis, e
margens materiais declaradas como não elicitadas --- não há, nesta unidade,
um limiar acordado a partir do qual um desvio importa. Nenhuma linha desta
secção se descreve como \emph{holdout}.

A classe de candidatos --- três estimadores, quatro construções de banda ---
e o critério que os adjudica estão definidos na
Secção~\ref{sec:protocolo-baselines}, com o que a aceitação de um candidato
estabelece e o que não estabelece; é dessa assimetria que saem as duas pistas
da Secção~\ref{sec:implicacoes-plataforma}. Esta secção apresenta o que a
aplicação desse critério decidiu.

\subsection{Estimadores robustos}
\label{sec:demo-robustos}

Para cada vetor e par anual estimaram-se três retas sobre o ano de treino
--- mínimos quadrados (o método de produção), Huber e Theil--Sen --- e
compararam-se as perdas absolutas diárias no ano seguinte, sem reestimação,
por contraste emparelhado $|e_{\mathrm{OLS}}| - |e_{\mathrm{rob}}|$ com
intervalo de confiança a 95\%. São quarenta contrastes.

\begin{table}[htbp]
  \centering
  \caption[Estimadores robustos contra mínimos quadrados]{Contrastes
  emparelhados de perda absoluta no ano seguinte, quatro vetores $\times$
  cinco pares. ``Robusto melhor'' e ``OLS melhor'': o intervalo MBB a 95\%
  exclui zero do lado respetivo.}
  \label{tab:robustos}
  \begin{tabular}{lccc}
    \toprule
    Estimador & Robusto melhor & OLS melhor & Sem evidência \\
    \midrule
    Huber ($c = 1{,}345$) & 11 & 6 & 3 \\
    Theil--Sen            & 7  & 8 & 5 \\
    \bottomrule
  \end{tabular}
\end{table}

Nenhum estimador domina (Tabela~\ref{tab:robustos}). O Huber ganha em onze
dos vinte contrastes e perde em seis; o Theil--Sen é mais extremo, sete
contra oito. O maior ganho é do Theil--Sen no vapor de 3~bar em
2022${\to}$2023 ($+4{,}68$~t~GNE/d, $+15{,}3\%$). Três leituras
qualificam a contagem. Contar vitórias não é agregar: no vapor de 24~bar o
Huber ganha três pares em cinco e tem, ainda assim, o erro absoluto médio
agregado ligeiramente pior. A explicação intuitiva --- ``a robustez ajuda
onde há caudas pesadas'' --- é contradita pelos próprios dados: o vapor de
3~bar, onde o Huber mais ganha, é o vetor com menor curtose residual, e o
fuel gás de 2024, com curtose $15{,}7$, não beneficia. E melhor ponto não é
melhor banda: no vapor de 24~bar em 2024${\to}$2025 o Huber reduz o erro
médio em $2{,}9\%$ e a banda marginal construída sobre a sua escala robusta
cobre $48{,}5\%$ dos dias, contra $93{,}9\%$ da banda de produção --- a
escala robusta é $0{,}55$ da escala dos mínimos quadrados, e uma banda mais
estreita em torno de um ponto marginalmente melhor é uma banda muito pior:
a nitidez só vale sujeita à calibração~\cite{Gneiting2007}.
A robustez na estimação e a robustez do detetor são problemas distintos, e
esta é a razão de a linha orientadora os separar.

\subsection{Bandas condicionais e quantílicas}
\label{sec:demo-bandas}

Das quatro construções definidas na Secção~\ref{sec:protocolo-baselines}, a
rota primária são as duas autorregressivas --- ordem entre um e sete dias,
predição a um passo no ano seguinte, com recurso declarado à banda marginal
quando o histórico não chega (124 dos $6\,580$ dias avaliados nas duas
construções, $1{,}9\%$), a gaussiana primária no veredicto --- e a rota de
sensibilidade é a regressão quantílica. Uma banda é
dada como suportada num vetor se o intervalo de confiança da sua cobertura
inclui $95{,}45\%$ em pelo menos quatro dos cinco pares, sem degeneração ---
compatibilidade pelo critério adotado, no sentido e com os limites acima.

\begin{table}[htbp]
  \centering
  \caption[Bandas alternativas nos vinte pares]{Cobertura no ano seguinte
  das quatro construções de banda, vinte pares vetor--transição. ``Inclui'':
  o intervalo MBB da cobertura inclui a referência de $95{,}45\%$; ``sub'' e
  ``sobre'': exclui-a por baixo ou por cima; NE: não estimável (não
  convergência declarada). Meia-largura mediana relativa à banda de
  produção.}
  \label{tab:bandas}
  \begin{tabular}{lccccc}
    \toprule
    Banda & Inclui & Sub & Sobre & NE & Meia-largura \\
    \midrule
    Produção $\pm 2\hat\sigma$   & 10 & 7  & 3 & 0 & $1{,}00$ \\
    AR, gaussiana                & 9  & 6  & 5 & 0 & $0{,}63$ \\
    AR, quantis empíricos        & 10 & 6  & 4 & 0 & $0{,}59$ \\
    Quantílica anual             & 5  & 12 & 2 & 1 & $0{,}97$ \\
    \bottomrule
  \end{tabular}
\end{table}

O veredicto é que \textbf{nenhuma rota candidata é suportada em nenhum dos
quatro vetores} (Tabela~\ref{tab:bandas}) --- e é preciso dizê-lo com
exatidão, porque a frase ``nenhuma banda é suportada'' seria falsa: a banda
de produção do vapor de 24~bar é compatível com a referência nos cinco pares.
O que a rota autorregressiva faz é instrutivo nas duas direções. Recupera
as falhas mais graves da produção --- no vapor de 10~bar em 2024${\to}$2025
a cobertura passa de $52{,}9\%$ para $95{,}6\%$ (mais $42{,}7$ pontos,
intervalo $[30{,}2;\,56{,}3]$); no de 3~bar em 2021${\to}$2022, de $37{,}8\%$
para $76{,}1\%$ --- e quebra três dos cinco sucessos do 24~bar, porque a sua
banda, $0{,}63$ da largura da produção, é estreita de mais onde a de
produção já era adequada. O saldo nos vinte pares é de menos um. Dois
pormenores mudam a leitura. A rota autorregressiva não é só uma banda: ao
usar os resíduos dos dias anteriores, muda o \emph{centro} da predição ---
é um modelo diferente, com memória, e parte do que recupera é o nível que os
mínimos quadrados anuais não seguem (o que a Secção~\ref{sec:referencia-congelada}
retoma). E a regressão quantílica sobre um ano tem sete ou oito observações
por cauda: a linha orientadora prescreve uma janela plurianual, a
demonstração executou a versão anual, e o seu resultado negativo é uma
sensibilidade dessa implementação, não um teste da recomendação
plurianual. As larguras da Tabela~\ref{tab:bandas} comparam-se na mesma
população de dias e ao mesmo nível nominal.

\subsection{O critério de aceitação e o que decidiu}
\label{sec:demo-gate}
A Figura~\ref{fig:matriz-decisao} mostra a matriz completa: cada candidato, cada par anual e o estado de cada um dos dois critérios.

\begin{figure}[p]
  \centering
  \includegraphics[width=\textwidth]{gerado/cap7-matriz-decisao}
  \caption[Matriz de decisão por candidato e por par anual]%
  {Os trinta e dois candidatos efetivamente avaliados --- quatro vetores, três ajustes pontuais e as bandas construídas sobre cada um --- contra os cinco pares anuais. Cada célula tem duas metades: à esquerda o critério de ganho preditivo, que é uma propriedade do ajuste pontual e por isso se repete entre bandas do mesmo estimador; à direita o critério de suporte da banda. As contagens à direita são os pares em que cada critério é cumprido, e os dois quadrados marcam se atingem a regra declarada de quatro em cinco. Nenhum candidato a cumpre nos dois critérios, e a regra de sensibilidade de três em cinco não altera essa conclusão. O conjunto avaliado não é o fatorial de estimadores por bandas: as bandas robustas só existem para os estimadores robustos.}
  \label{fig:matriz-decisao}
\end{figure}

O critério foi aplicado por adjudicação pura sobre os resultados já
publicados --- o desempenho preditivo do Capítulo~\ref{cha:diagnostico} e as
bandas da secção anterior ---, sem recalcular nada, com cada valor lido de
uma corrida identificada e verificada por \emph{hash}. Com os mínimos quadrados e as quatro construções de banda, são dezasseis
decisões. O par junta dois objetos julgados em separado: o normalizador, a reta anual, pelo conteúdo preditivo do
Capítulo~\ref{cha:diagnostico}, e a banda, pela cobertura. Nas construções
autorregressivas a banda desloca também o centro da predição com os
resíduos dos dias anteriores; o erro dessa predição completa não foi
adjudicado, e o zero em dezasseis refere-se a esta regra, não à inutilidade
da predição autorregressiva.

\begin{table}[htbp]
  \centering
  \caption[Critério de aceitação: candidatos de mínimos quadrados]{Adjudicação
  dos dezasseis candidatos com estimador de mínimos quadrados. Conteúdo
  preditivo: pares (em cinco) com ganho sobre a média do treino e intervalo a
  excluir zero; banda: pares com cobertura compatível com a nominal. Aceite
  se ambos $\geq 4/5$.}
  \label{tab:gate}
  \small
  \begin{tabular}{lcccccc}
    \toprule
    Vetor & Cont. preditivo & Produção & AR gauss. & AR emp. & Quantílica & Aceites \\
    \midrule
    Fuel gás     & \textbf{4/5} & 3/5 & 3/5 & 3/5 & 2/5 + NE & 0/4 \\
    Vapor 24 bar & 1/5 & \textbf{5/5} & 2/5 & 2/5 & 2/5 & 0/4 \\
    Vapor 10 bar & 0/5 & 1/5 & 3/5 & 3/5 & 0/5 & 0/4 \\
    Vapor 3 bar  & 2/5 & 1/5 & 1/5 & 2/5 & 1/5 & 0/4 \\
    \bottomrule
  \end{tabular}
\end{table}

\textbf{Nenhum dos dezasseis candidatos é aceite} (Tabela~\ref{tab:gate}), e
as razões separam-se limpamente por vetor. O fuel gás cai só pela banda:
tem conteúdo preditivo em quatro dos cinco pares, mas nenhuma construção o
acompanha com cobertura compatível em mais de três. Os vapores caem pelo
conteúdo preditivo --- um, zero e dois pares em cinco --- e a única célula
em que uma banda passa, a de produção do 24~bar, \revisto{acompanha uma reta sem ganho preditivo consistente sobre a média. O contraste com o fuel gás separa os dois critérios, sem retirar à média o estatuto de modelo.}

Os candidatos robustos completos --- Huber e Theil--Sen, cada um com a banda
marginal sobre a sua escala robusta e com a rota autorregressiva sobre os
seus resíduos, mais dezasseis decisões --- não mudam o desfecho: zero
aceites. Os quatro candidatos do fuel gás mantêm o conteúdo preditivo
(quatro em cinco com ambos os estimadores) e nenhuma banda passa de três em
cinco; os vapores continuam sem conteúdo preditivo (o melhor, Theil--Sen no
3~bar, chega a três em cinco). Trocar o estimador pontual não resgata a
classe. A correção autorregressiva sobe o número de pares compatíveis de
nove, com a banda marginal, para dezoito --- sem nunca chegar a quatro num
mesmo candidato; o caso extremo é o vapor de 24~bar com Theil--Sen em
2024${\to}$2025, em que a banda marginal cobre $14{,}6\%$ e a autorregressiva
$96{,}6\%$.

O que o zero em trinta e dois significa, e o que não significa, tem de ficar
escrito com o mesmo cuidado. Significa que nenhuma das combinações avaliadas
--- três estimadores, seis construções de banda, quatro vetores --- cumpre
o critério declarado sob a regra de quatro em cinco. Não demonstra a
impossibilidade da classe linear, nem a necessidade de modelos com estado
--- essa frase chegou a constar das fichas e foi retirada. E o resultado é
sensível ao limiar: \revisto{com uma regra de três em cinco, quatro dos dezasseis candidatos de mínimos quadrados passariam o critério de conteúdo preditivo e seis dos dezasseis robustos, mas nenhum dos dez seria aceite, por nenhum ter banda suportada;} a regra de quatro em cinco estava pré-registada e não
se altera depois de conhecido o resultado, mas o leitor deve saber que a
decisão assenta em cinco pares anuais, e que a potência de uma regra binária
sobre cinco observações dependentes nunca foi calculada (Secção~\ref{sec:limitacoes}).
O que a demonstração estabelece de sólido é o \emph{procedimento}: um
critério declarado a priori, adjudicação separada do cálculo, razões de
rejeição distintas por vetor, e um resultado negativo que fica publicado
como tal em vez de ser ajustado até aparecer um vencedor.

\subsection{Fichas por \emph{baseline}}
\label{sec:demo-fichas}

A última linha orientadora foi demonstrada produzindo, para cada um dos
quatro vetores, a ficha de uma página que a governança exige: domínio de
validade, janela, exclusões, estimador, método de banda e o seu desfecho,
resultados do critério de aceitação, e proveniência --- setenta e seis
campos por vetor, cada um lido de uma corrida identificada e verificado por
\emph{hash} contra o manifesto de origem, sem nenhum valor escrito à mão. As
fichas não calculam nada: consolidam. E é exatamente por isso que servem de
especificação para a plataforma (Secção~\ref{sec:implicacoes-plataforma}): o
que uma pessoa precisa de saber para confiar numa \emph{baseline} cabe numa
página, desde que cada campo tenha origem rastreável. Duas erratas
apanhadas em revisão --- um máximo de 2026 publicado como se fosse do
período 2020--2025, e uma frase sobre ``deriva de nível como falha estrutural
documentada'' que as corridas não sustentam --- foram corrigidas na versão
final.

\section{Referência congelada e verdade conhecida}
\label{sec:referencia-congelada}

A demonstração anterior manteve a perspetiva do método de produção:
reajustar todos os anos e perguntar se o ajuste transporta. A exploração
que se segue muda de perspetiva. Congela uma \emph{referência} --- a relação
estimada em 2020, com a sua escala de erro e, quando aplicável, o seu estado
inicial --- e aplica-a a 2021--2025 sem lhe tocar. Tudo o que se mede na
aplicação passa a ser desvio face a 2020, e o desenho separa três objetos
que o método de produção mistura num só: o resíduo da referência (o desvio
físico face ao ponto fixo), o estado filtrado (o que um modelo adaptativo
segue) e a inovação (o que o filtro não previu). O reajuste anual é aqui
uma sensibilidade, não o primário. Todos os resumos usam uma população
declarada --- dias de 2021--2025 retidos pela seleção oficial, dentro do
suporte de carga de 2020, com observação e predição finitas --- e uma
comparação entre modelos faz-se sempre na interseção dessas populações; foi
a imposição desta regra, numa revisão independente, que inverteu várias
ordenações da primeira versão desta exploração, e a versão que aqui se cita
é a corrigida.

\subsection{O que muda quando não se reajusta}
\label{sec:referencia-fixa}

\paragraph{A referência congelada mede o desvio.} \ajuste{A escala da
referência é fixada por exclusão de blocos contíguos
(Equação~\ref{eq:mat-crossfit}), e o que um reajuste faz ao sinal está na
Equação~\ref{eq:mat-projecao-sinal}.} Aplicada a 2021--2025 na
população comum aos cinco normalizadores avaliados ($n = 876$ dias), a
relação de 2020 tem, no fuel gás, \emph{skill} de $41$ a $54\%$ face à média
de 2020 --- os normalizadores continuam a bater a média cinco anos depois. A
amostra comum são 876 dos $1\,645$ dias retidos de 2021 a 2025 em cada vetor
($53\%$): exigir carga dentro do suporte de 2020 deixa $1\,232$, e exigir
densidade e condições de \emph{flash} dentro do intervalo de 2020, como o
normalizador com covariáveis impõe, deixa os 876. Comparar na interseção é
justo entre modelos, mas pode favorecer os períodos mais parecidos com 2020;
na população própria de cada modelo, os valores mudam pouco (fuel gás de
$42$ a $54\%$; vapor de 24~bar de $-6$ a $0\%$; 10~bar de $2$ a $6\%$;
3~bar de $-15$ a $-23\%$). O ano de 2020 é a referência por ser o primeiro
ano completo disponível, não por ser um regime comprovadamente limpo ou
eficiente.
No vapor de 24~bar nenhum normalizador acrescenta nada à média de 2020 (de
$-5{,}6\%$ a $+0{,}3\%$); no de 10~bar o ganho é de $2$ a $6\%$; no de
3~bar a média de 2020 é o melhor ``modelo'' de todos, por $16$ a $25\%$. O
desvio padronizado médio face à referência é de $0{,}97\sigma$ no fuel gás
e $1{,}73\sigma$ no vapor de 10~bar: a aplicação já não se parece com o ano
de referência, e um normalizador congelado em 2020 é o que mede essa
distância. Uma banda de $95\%$ calibrada em 2020 cobre, em 2021--2025 e nos
dias em que a comparação é admissível, entre $0{,}564$ e $0{,}984$ dos dias
consoante o vetor, o modelo e a regra --- mediana entre vetores de
$0{,}71$ a $0{,}86$ por modelo ---, com o vapor de 24~bar acima de $0{,}94$
em todos os tercis de carga e o de 10~bar a cair a $0{,}43$ no tercil de
carga baixa. Isto é a medição, não um defeito a corrigir: a banda diz que a
aplicação já não se parece com a referência, e nenhuma banda é alargada até
o número voltar. Perder cobertura nominal face a uma referência congelada
não é, por si, um veredicto de ineficiência (Secção~\ref{sec:limitacoes});
é a quantidade que um sistema de monitorização deve mostrar e o método de
produção reabsorve.

\paragraph{Um controlo adaptativo cobre mais, por seguir o desvio.} Nos
mesmos dias, um controlo por quantis móveis dos noventa dias civis
anteriores --- só com passado admissível --- cobre $0{,}87$ a $0{,}91$ dos
dias nos quatro vetores, contra $0{,}60$ a $0{,}99$ da banda fixa; e fá-lo
com uma banda mais estreita, em unidades da escala de 2020, no vapor de
10~bar ($3{,}02$ contra $4{,}48$) e no de 24~bar, e mais larga no fuel gás.
Cobre mais porque se adapta ao desvio que a banda fixa existe para revelar
--- é a Secção~\ref{sec:leitura-integrada} em números --- e por isso não é
um substituto: é um controlo, pontuável em $38\%$ dos dias primários
(exige sessenta erros admissíveis nos noventa dias anteriores).

A Figura~\ref{fig:injeccao} mostra, dia a dia, quanto do desvio inserido cada componente ainda deixa ver.

\begin{figure}[tbp]
  \centering
  \includegraphics[width=\textwidth]{gerado/cap7-injeccao-degrau}
  \caption[Quanto sinal sobrevive a cada componente]%
  {Experiência de injeção no fuel gás: um desvio conhecido é somado à série e segue-se o que resta dele em cada componente, nos primeiros 150 dias de aplicação, para um degrau e para uma rampa da mesma amplitude. Em cima, o desvio inserido. Ao meio, o desvio medido contra a referência congelada e a parte que o reajuste anual absorve. Em baixo, a inovação do filtro. A diferença entre os dois primeiros painéis é uma identidade algébrica, não uma medição; o terceiro painel é que mostra a atenuação efetiva. Os valores numéricos citados no texto vêm do lote validado, não desta ilustração.}
  \label{fig:injeccao}
\end{figure}
\paragraph{Um desvio inserido sobrevive à referência fixa e não sobrevive ao
reajuste.} Para medir o mecanismo diretamente, inseriu-se um desvio
conhecido --- um degrau e uma rampa, em unidades da escala de cada vetor
--- sobre os dados reais de aplicação, e mediu-se a fração desse desvio que
sobrevive em cada canal. O resíduo da referência congelada preserva-o por
inteiro, em qualquer das três leituras (janela inteira, dias em amplitude
plena, primeira quinzena em amplitude plena) --- é uma identidade algébrica,
não um mérito, e é o que ``referência'' significa. Um reajuste por mínimos quadrados
sobre o próprio período desviado absorve na ordenada na origem um degrau
constante nesse período e deixa no resíduo $3\%$ de uma rampa. É a parte
algébrica do mecanismo, não uma impossibilidade de deteção: uma mudança que
afete só parte da janela de reajuste, ou os dados posteriores, não
desaparece assim. E a inovação de um filtro de nível~\cite{DurbinKoopman2012} assimila um
desvio sustentado em poucos dias: da primeira quinzena de um degrau retém
$0{,}136$; da janela inteira, $0{,}002$; uma rampa nunca chega a produzir
inovação, porque sobe mais devagar do que o filtro aprende. Esta última
leitura é sobre a inovação de um filtro, não sobre uma estatística de
monitorização --- se um monitor deteta a mudança é a pergunta da secção
seguinte. O que fica medido, sob este desenho, é a atenuação do
mecanismo descrito no Capítulo~\ref{cha:diagnostico}\matref{, quantificada
pela Equação~\ref{eq:mat-atenuacao-ar}}: sobre o período usado
para reajustar, o reajuste absorve o degrau e quase toda a rampa.

\paragraph{Janelas de reajuste: hipótese pendente.} A sensibilidade que
reajusta a referência com janelas móveis de seis, doze e vinte e quatro
meses, ou expansiva~\cite{PesaranTimmermann2007}, está entre as partes em
revisão (fim desta secção): compara populações diferentes, e nenhuma
ordenação das janelas se usa nesta dissertação para justificar uma
recomendação. Fica como hipótese a testar depois da correção que a janela
mais curta siga melhor onde há desvio a seguir e, por isso mesmo, o
acomode mais depressa.

\subsection{Monitor e \emph{baseline} com verdade conhecida}
\label{sec:verdade-conhecida}

A pergunta que o calendário operacional tornou não estimável no
Capítulo~\ref{cha:diagnostico} --- que desvio o sistema deteta, e o que
importa mais para o detetar --- só pode responder-se com verdade conhecida.
Construiu-se, para cada vetor, um gerador \ajuste{(Equações~\ref{eq:mat-dgp}
a~\ref{eq:mat-rampa})} com a relação de 2020 (mínimos quadrados com as
covariáveis de referência), ruído autorregressivo de primeira ordem,
\matref{\begin{equation}
  \varepsilon_t = \rho\,\varepsilon_{t-1} + \sigma\sqrt{1-\rho^2}\,z_t ,
  \label{eq:gerador-ruido}
\end{equation}}
com $\rho = 0{,}6$ e $z_t$ normal padrão, de modo que a dispersão marginal
$\sigma$ é a dos resíduos de 2020 e a das inovações é $\sigma\sqrt{1-\rho^2}$,
366 dias de treino com a máscara real de dias admissíveis de 2020 e uma
aplicação de noventa dias. As alternativas medem-se em unidades de $\sigma$
marginal e começam no primeiro dia da aplicação (no dia 45, em
sensibilidade); o canal de nível é unilateral superior no fuel gás e nos
vapores de 24 e 3~bar e bilateral no de 10~bar, que é um fluxo líquido. O
limiar de cada canal, em cada política, é o quantil de 95\% do máximo da
sua estatística nos noventa dias numa coorte de réplicas de calibração
\ajuste{(Equação~\ref{eq:mat-fap-horizonte})}: a
taxa de falso sinal de $5\%$ é a probabilidade de pelo menos um sinal no
horizonte, por canal. Foi verificada numa coorte independente (quarenta
canais entre $0{,}020$ e $0{,}072$, um com o intervalo de
Wilson~\cite{Wilson1927} inteiramente acima de $0{,}05$; com quarenta canais
seria de esperar um assim por acaso, mas não havia um critério global
declarado para o afirmar), e as alternativas usam uma terceira coorte. Os
canais avaliam-se separadamente; o único canal conjunto é o do fim desta
subsecção, com calibração própria. Sobre este gerador inseriram-se
alternativas conhecidas --- degrau e rampa de uma escala, alargamento da
dispersão, recuperação, resposta à carga --- em quatrocentas réplicas por
célula, e mediu-se a taxa de deteção de cada combinação de \emph{baseline}
(mínimos quadrados ou modelo
aditivo~\cite{HastieTibshirani1986,Wood2017}\matref{; Equações~\ref{eq:mat-gam}
e~\ref{eq:mat-penalizacao}}) e
monitor (resíduo bruto, filtro de
nível, soma acumulada, e canais de carga e de
escala)\matref{~---~o modelo de nível local e o seu passo de atualização estão
nas Equações~\ref{eq:mat-nivel-local} a~\ref{eq:mat-kalman-variancia}, a soma
acumulada na Equação~\ref{eq:mat-cusum} e os dois canais nas
Equações~\ref{eq:mat-monitor-carga} e~\ref{eq:mat-monitor-escala}~---}. \textbf{Nada nesta
secção é um facto sobre a refinaria}: é o comportamento de métodos sobre um
gerador que escrevemos, e deve ser lido como tal.

\begin{table}[htbp]
  \centering
  \caption[Contrastes co-primários com verdade conhecida]{Diferença de taxa
  de deteção de um degrau de $1\sigma$, emparelhada por réplica, com intervalo
  percentil de 2\,000 reamostragens. Esquerda: trocar a \emph{baseline}
  (mínimos quadrados $\to$ modelo aditivo) com o monitor de nível fixo.
  Direita: trocar o monitor (filtro de nível $\to$ soma acumulada) com a
  \emph{baseline} fixa. Gerador com $\rho = 0{,}6$.}
  \label{tab:coprimarias}
  \begin{tabular}{lcccc}
    \toprule
    Vetor & Trocar a \emph{baseline} & IC 95\% & Trocar o monitor & IC 95\% \\
    \midrule
    Fuel gás     & $-0{,}028$ & $[-0{,}052;\,-0{,}005]$ & $+0{,}323$ & $[0{,}273;\,0{,}370]$ \\
    Vapor 24 bar & $-0{,}013$ & $[-0{,}040;\,0{,}015]$  & $+0{,}438$ & $[0{,}388;\,0{,}487]$ \\
    Vapor 10 bar & $-0{,}020$ & $[-0{,}045;\,0{,}003]$  & $+0{,}490$ & $[0{,}440;\,0{,}540]$ \\
    Vapor 3 bar  & $-0{,}028$ & $[-0{,}060;\,0{,}003]$  & $+0{,}390$ & $[0{,}343;\,0{,}438]$ \\
    \bottomrule
  \end{tabular}
\end{table}

\paragraph{Sob o gerador linear, o monitor pesa mais do que a
\emph{baseline}.} As oito comparações declaradas como primárias antes da
execução (Tabela~\ref{tab:coprimarias}) dizem, nos quatro vetores, a mesma
coisa: trocar o monitor de nível por uma soma acumulada aumenta a taxa de deteção
de um degrau de uma escala em 32 a 49 pontos percentuais, com os quatro
intervalos longe de zero; trocar a \emph{baseline} linear por um modelo
aditivo muda-a em um a três pontos, e só no fuel gás o intervalo exclui zero
--- a favor da linear. Cada comparação vale por si: as oito não têm correção
de multiplicidade. E o gerador primário é linear, o que favorece a
\emph{baseline} linear; um ganho pequeno do modelo aditivo neste gerador não
mostra que a escolha da \emph{baseline} importe pouco em geral. Os quatrocentos e quarenta contrastes secundários,
publicados sem nunca serem promovidos, dizem onde isto \emph{não} vale: num
alargamento puro de dispersão a soma acumulada \emph{perde} ($-0{,}28$ em
média), porque acumula desvios de média e esses cancelam-se; em rampas e
recuperações compra $0{,}29$ e $0{,}20$; e num banco com interação
carga$\times$densidade a \emph{baseline} aditiva ganha. A escolha do monitor
é um compromisso, não uma melhoria --- e nenhuma recomendação vale sem a
natureza do desvio e a autocorrelação medidas ao lado.

A Figura~\ref{fig:deteccao} compara os três monitores sob o mesmo orçamento de falsos sinais.

\begin{figure}[tbp]
  \centering
  \includegraphics[width=\textwidth]{gerado/cap7-deteccao-simulacao}
  \caption[Deteção simulada, ao mesmo orçamento de falsos sinais]%
  {Probabilidade de detetar um desvio nos primeiros 90 dias, por amplitude, para três monitores calibrados ao mesmo orçamento de falsos sinais. Os dois primeiros painéis separam degrau de rampa com dependência moderada ($\rho=0{,}6$); o terceiro fixa a amplitude num desvio-padrão e varia a dependência do gerador. A banda é o intervalo de Wilson sobre as réplicas. A linha a tracejado marca metade das deteções. Os valores descrevem o gerador simulado e as suas hipóteses, e não são uma previsão do que aconteceria na refinaria.}
  \label{fig:deteccao}
\end{figure}
\matref{A dependência entre vetores e o efeito da agregação temporal sobre a
variância estão nas Equações~\ref{eq:mat-ar-multivariado}
e~\ref{eq:mat-variancia-soma}.}

\paragraph{A dependência decide quanto se deteta.} O gerador primário tem
$\rho = 0{,}6$, que é a autocorrelação residual observada no fuel gás; os
três vapores estão entre $0{,}8$ e $0{,}9$, e a subir
(Secção~\ref{sec:diagnosticos-desc}). O lote de sensibilidade responde ao que
isso muda. Num degrau de uma escala, a soma acumulada deteta $0{,}96$ das
réplicas a $\rho = 0{,}6$, $0{,}59$ a $\rho = 0{,}9$ e $0{,}39$ a
$\rho = 0{,}95$ --- neste último caso igual ao resíduo bruto ($0{,}35$): a
vantagem do monitor desaparece. Numa rampa, a taxa cai de $0{,}64$ para
$0{,}27$ e $0{,}17$. A leitura para os vapores é, por isso, mais modesta do
que a das oito células primárias: na dependência que eles exibem, um degrau
de uma escala é detetado dentro do horizonte em cerca de metade das
réplicas pelo melhor monitor, e uma rampa em um quarto. O monitor continua
a pesar; a detetabilidade absoluta é pobre, e o resultado do fuel gás não
transfere.

\paragraph{O desvio mínimo detetável, no gerador.} Sobre a curva de deteção
por amplitude de degrau \ajuste{(Equações~\ref{eq:mat-deteccao}
a~\ref{eq:mat-mdd-interpolacao})}, o desvio a que a soma acumulada deteta metade das
vezes situa-se entre $0{,}51\sigma$ (vapor de 3~bar) e $0{,}63\sigma$ (vapor
de 10~bar), contra $0{,}87$ a $1{,}10\sigma$ do filtro de nível e $0{,}95$ a
$1{,}15\sigma$ do resíduo bruto. Em percentagem do consumo médio dos dias
retidos de 2020, o desvio mínimo detetável da soma acumulada é $4{,}5\%$ no
fuel gás, $9{,}4\%$ no vapor de 3~bar e $14{,}7\%$ no de 24~bar; o vapor de
10~bar é um fluxo líquido com sinal e não admite percentagem. Trinta e duas
das oitenta curvas estão censuradas à direita --- todas nos canais de carga e
de escala, que não existem para ver um degrau de nível --- e são publicadas
como limite aberto, nunca como número. \revisto{Estes valores são propriedades do gerador a $\rho = 0{,}6$; o parágrafo anterior mostra a sua sensibilidade à dependência, sem estimar a detetabilidade real na refinaria.}

\paragraph{Três sinais falsos, demonstrados.} O mesmo aparelho serviu para
demonstrar, com verdade conhecida, três leituras sedutoras que a prática de
monitorização produz. Primeira: numa realização, duas séries \emph{independentes} têm
$R^2$ diário de $0{,}0025$, mensal de $0{,}18$ e acumulado de $0{,}9999$ ---
a correlação de séries acumuladas pode ser fabricada pela
acumulação~\cite{Yule1926,GrangerNewbold1974}, e um gráfico de consumo
acumulado contra carga acumulada não é, por si, evidência de relação. A
distribuição destes valores por réplicas não foi tabulada, e um $R^2$
acumulado alto não torna espúria toda a relação acumulada. Segunda: branquear a série antes de monitorizar --- usar as
inovações de um modelo autorregressivo como estatística --- atenua a mudança
de nível que se quer ver~\cite{AlwanRoberts1988,Wardell1994}, porque o
modelo passa a prever o degrau a partir dos próprios desfasamentos e, num
processo autorregressivo de primeira ordem, a inovação só retém a fração
$1-\rho$ de um degrau sustentado. Num degrau de uma escala, com cada
estatística calibrada para $5\%$ de falsos sinais, a inovação
autorregressiva sinaliza em $0{,}095$ das réplicas, pouco acima do que
sinalizaria sem desvio, enquanto uma média móvel de trinta dias deteta
$0{,}895$. Terceira: a leitura de que a
anticorrelação entre os vapores de 24 e 10~bar
(Secção~\ref{sec:diagnosticos-desc}) ``impede a coocorrência'' de alarmes ---
que constou de uma versão anterior desta exploração e foi retirada ---
assentava em pontuar o vapor de 10~bar como unilateral; na lateralidade
declarada, com os quatro vetores desviados, um canal conjunto sinaliza em
$0{,}068$ das réplicas contra $0{,}570$ da melhor margem individual, e com um
só vetor desviado em $0{,}022$. Uma contagem conjunta exige mais do que a
melhor margem e por isso deteta menos; não é a correlação que o impede.

\paragraph{O que a exploração com referência congelada ainda não fecha.}
Esta exploração foi submetida a duas revisões independentes. A primeira
aprovou a integridade da cadeia de resultados e reprovou a conformidade
científica de sete pontos, todos corrigidos e republicados. A segunda,
sobre a versão republicada, identificou que várias das suas partes --- a
comparação entre esquemas de referência, a ordenação de janelas de
reajuste, a inserção de desvios sobre estados publicados, a agregação a
frequências semanal e mensal e o suporte da eletricidade --- ainda não
estimam exatamente a comparação declarada (comparam populações diferentes,
comprimem o calendário, ou incluem extrapolações no resultado primário). As
correções estão em curso à data desta versão; este capítulo cita, com
números, apenas os lotes que essa revisão classificou como aceitáveis, e
onde uma parte em revisão entra, entra qualitativamente e assinalada.

\section{Implicações para a plataforma}
\label{sec:implicacoes-plataforma}
A Figura~\ref{fig:ficha-estados} concretiza os requisitos numa ficha e num ciclo de vida.

\begin{figure}[htbp]
  \centering
  \textbf{\small (a) Ficha da referência --- proposta, não implementada}\\[3pt]
  \begin{tikzpicture}[x=1cm, y=1cm, font=\scriptsize,
      cmp/.style={draw=corCinzaC, line width=.5pt, fill=white, align=left,
                  inner sep=3pt, font=\scriptsize, text=corTinta,
                  text width=3.9cm, minimum height=13mm, anchor=north west}]
    \providecommand{\rot}{}\renewcommand{\rot}[1]{{\tiny\color{corCinza}#1}}
    \begin{scope}[on background layer]
      \fill[corFundo, rounded corners=2pt] (0,-3.95) rectangle (13.6,1.30);
      \draw[corCinza, line width=.6pt, dash pattern=on 3pt off 2pt,
            rounded corners=2pt] (0,-3.95) rectangle (13.6,1.30);
    \end{scope}
    \node[anchor=north west, font=\scriptsize\bfseries, text=corTinta] at (0.30,1.15)
      {Destilação atmosférica --- fuel gás};
    \node[anchor=north east, font=\tiny, text=corCinza] at (13.30,1.12)
      {referência v4, de 2026-03-02 \quad$\vert$\quad versão anterior consultável};
    \draw[corCinzaC, line width=.5pt] (0.30,0.62) -- (13.30,0.62);

    \node[cmp] at (0.30,0.42)
      {\rot{DOMÍNIO DE VALIDADE}\\ carga entre 18,6 e 27,4~kt/d\\
       \rot{fora do domínio não há juízo}};
    \node[cmp] at (4.80,0.42)
      {\rot{DESVIO DO DIA}\\ $+$41 t GNE \quad($+$12,8\,\%)\\
       \rot{observado $-$ esperado}};
    \node[cmp] at (9.30,0.42)
      {\rot{INCERTEZA}\\ banda de 95\,\%: $\pm$94 t GNE\\
       \rot{cobertura verificada: 0,87}};

    \node[cmp, draw=corFG, line width=.7pt, fill=corFG!7] at (0.30,-1.20)
      {\rot{ESTADO}\\ \textbf{em observação}\\
       \rot{desde 2026-07-14, 4.\textordmasculine{} dia}};
    \node[cmp] at (4.80,-1.20)
      {\rot{AÇÃO EM CURSO}\\ verificar medição do canal\\
       \rot{responsável: engenharia de processo}};
    \node[cmp] at (9.30,-1.20)
      {\rot{QUALIDADE DO DIA}\\ 3 origens concordantes\\
       \rot{sem imputação}};

    \draw[corCinzaC, line width=.5pt] (0.30,-2.85) -- (13.30,-2.85);
    \node[anchor=north west, font=\tiny, text=corCinza, text width=13.0cm]
      at (0.30,-2.98)
      {Seis campos na ficha principal. Os restantes campos do registo ---
       parâmetros, histórico de versões, diagnóstico completo --- ficam no
       registo detalhado (Apêndice~\ref{app:manual}) e não competem com estes
       pela atenção de quem decide.};
  \end{tikzpicture}

  \vspace{4mm}
  \textbf{\small (b) Estados da referência e eventos que os mudam}\\[3pt]
  \begin{tikzpicture}[x=1cm, y=1cm, font=\scriptsize,
      st/.style={ntcaixa, text width=2.3cm, minimum height=10mm,
                 draw=corFG, line width=.7pt, fill=corFG!7},
      ev/.style={font=\tiny, text=corCinza, align=center, text width=1.45cm,
                 inner sep=1.5pt, fill=white}]
    \node[st] (op) at (1.45,0) {\textbf{em vigor}};
    \node[st] (ob) at (5.35,0) {\textbf{em observação}};
    \node[st] (su) at (9.25,0) {\textbf{suspensa}};
    \node[st] (ve) at (13.15,0) {\textbf{em verificação}};

    \draw[ntsetaf] (op) -- (ob);
    \node[ev, anchor=south] at (3.40,0.12) {sinal do monitor};
    \draw[ntsetaf] (ob) -- (su);
    \node[ev, anchor=south] at (7.30,0.12) {causa confirmada ou domínio violado};
    \draw[ntsetaf] (su) -- (ve);
    \node[ev, anchor=south] at (11.20,0.12) {reestimação datada};

    \draw[ntsetaf] (ve.north) to[out=105, in=75, looseness=0.55] (op.north);
    \node[ev, anchor=south, text width=5.0cm] at (7.30,1.55)
      {cobertura verificada fora da amostra};
    \draw[ntsetaf] (ob.south) to[out=-110, in=-70, looseness=0.9] (op.south);
    \node[ev, anchor=north, text width=4.2cm] at (3.40,-1.34)
      {sinal não confirmado em 10 dias};

    \node[ntmini, anchor=north west, text width=13.5cm] at (0.20,-2.10)
      {Os estados persistem e ficam datados; os eventos que os mudam são
       instantâneos e ficam registados com quem os desencadeou. Uma
       referência suspensa não produz juízo: a unidade continua a ser
       medida, mas não é julgada até a nova referência passar a verificação.};
  \end{tikzpicture}
  \caption[Ficha da referência e diagrama de estados]%
  {Proposta de ficha da referência e do ciclo de vida que a governa. Os
   valores mostrados são ilustrativos. O painel~(a) mostra o que quem
   decide precisa de ver de uma vez: onde a referência é válida, qual o
   desvio, com que incerteza, em que estado está, o que se está a fazer e
   que qualidade tem o dado do dia. O painel~(b) separa aquilo que o método
   em produção não separa: o alerta é um evento, o estado da referência é
   uma condição persistente, e a passagem de um ao outro exige uma decisão
   com dono e data. Nada disto foi implementado neste trabalho.}
  \label{fig:ficha-estados}
\end{figure}


O Capítulo~\ref{cha:plataforma} entregou um instrumento que torna o
diagnóstico possível; este capítulo devolve-lhe uma especificação. O que se
segue é o bloco de saúde das \emph{baselines} --- o que o \emph{dashboard}
tem de mostrar para que uma \emph{baseline} seja usada com conhecimento do
seu estado --- formulado como requisitos, cada um ancorado num resultado
desta dissertação.

\paragraph{Duas pistas, não uma.} O método de produção usa um só modelo para
normalizar e para alarmar. A plataforma deve manter duas pistas visíveis
lado a lado: o modelo anual em vigor, para reporte normalizado e comparação
entre períodos, como a ISO~50006 pede; e uma referência congelada, com o
desvio padronizado face a ela e um monitor sequencial sobre esse desvio,
para deteção. \revisto{A Secção~\ref{sec:referencia-fixa} mostrou que a segunda conserva um termo de comparação fixo entre anos;} a Secção~\ref{sec:verdade-conhecida} mostrou que o monitor
sobre ela é o que mais pesa na deteção. A data e a razão de cada mudança de
referência são registadas.

\paragraph{Domínio visível, alarme suprimido fora dele.} Cada
\emph{baseline} publica o intervalo de carga e as condições em que foi
estimada; um dia fora desse domínio é mostrado como ``não avaliado'', não
como anomalia nem como operação normal, e o seu alarme é suprimido. O limiar de carga oficial já
define um domínio implícito (Secção~\ref{sec:perfil-dados}); a plataforma
torna-o explícito.

\paragraph{Calibração monitorizada, não assumida.} A cobertura efetiva da
banda de alarme nos últimos noventa e cento e oitenta dias, com intervalo de
confiança que respeite a dependência, é um indicador do bloco de saúde.
Cobertura e taxa de dias fora da banda são a mesma medição, uma o
complemento da outra, e nenhuma separa sozinha uma banda descalibrada de um
desvio real; o que acrescenta informação é o intervalo, a separação por
sentido (acima e abaixo) e por episódio (uma sequência de dias fora conta
como um episódio) e a investigação externa. Janelas móveis sobrepostas não
são testes independentes repetidos. Uma cobertura
que exclui a nominal pelo intervalo desencadeia a política de investigação,
não uma recalibração automática.

\paragraph{Ficha por \emph{baseline}.} Os setenta e seis campos da
Secção~\ref{sec:demo-fichas} --- domínio, janela, exclusões com códigos de
razão, estimador, método de banda e desfecho, resultados do critério de
aceitação, proveniência --- são a ficha que acompanha cada \emph{baseline}
na plataforma, gerada dos resultados e nunca editada à mão. Uma
\emph{baseline} sem ficha não está em produção.

\paragraph{Políticas separadas e registadas.} Alarme, investigação, suspensão,
recalibração e verificação pós-recalibração são cinco estados com transições
declaradas; a recalibração exige sinal do monitor mais revisão documentada,
com a opção ``causa desconhecida, recalibrar com sinalização''; a
certificação do modelo recalibrado usa só dados posteriores. O registo de
versões, aprovações e revisões é parte do bloco, não um documento à parte.

\paragraph{O que a plataforma não pode fazer sem informação que não tem.}
Nenhum dos requisitos acima classifica um desvio como ineficiência, nem um
alarme como falso. Isso exige o calendário operacional, rótulos de eventos e
margens materiais elicitadas com a operação --- as três coisas que esta
dissertação não teve. A plataforma pode e deve registar o desvio e o alarme
com a sua proveniência; a classificação é uma decisão humana até que essa
informação exista, e o desenho deve deixar espaço para a receber.

\section{Limitações}
\label{sec:limitacoes}

\paragraph{Do estudo.} Uma unidade, seis anos completos, um único
\emph{snapshot} de dados já explorados. Os pares são avaliados para a frente
no tempo, mas os mesmos anos serviram para formular as decisões, pelo que
nenhum resultado é uma validação independente da seleção: a demonstração das
linhas orientadoras é retrospetiva e encadeada, nunca \emph{holdout}, porque
2025 e o parcial de 2026 já tinham informado o diagnóstico. As conclusões são
condicionais ao modelo linear, à seleção oficial e às janelas anuais que o
método de produção herdou. O calendário operacional, os rótulos de eventos e
as margens materiais não estavam disponíveis: nenhum desvio é ineficiência,
nenhuma excursão é falso alarme, nenhuma cobertura é qualidade de deteção,
e a detetabilidade só foi estudada em simulação. Há erro de medição na
carga, possível endogeneidade (a operação reage aos alarmes), agregação
diária que esconde dinâmica intra-diária, um mecanismo de falta não
identificável sem fonte operacional, e nenhum padrão de referência
independente do modelo Power BI original para as observações selecionadas.
A energia elétrica é uma alocação mensal e fica fora de qualquer inferência
diária.

\paragraph{Dos dados.} A validação célula a célula cobre janeiro a maio de
2026 e o fecho de 400 em 400 é, nesta versão, esperado e não observado
(Secção~\ref{sec:validacao-dados}). Nos seis anos avaliados, o canal de
\emph{tags} e o canal de balanço da destilação atmosférica afastam-se de
forma que muda de ano para ano (Secção~\ref{sec:modelo-dimensional}); a
série avaliada é a do balanço, e não foi verificado que as mudanças entre
anos que o Capítulo~\ref{cha:diagnostico} encontra não reflitam também
alterações de medição ou de contabilização no próprio balanço.

\paragraph{Da inferência.} Todos os veredictos assentam em cinco pares
anuais e em regras binárias sobre eles; a potência dessas regras sob a
dependência observada não foi calculada, e com uma regra de três em cinco
três dos dezasseis candidatos de mínimos quadrados passariam. \revisto{O intervalo de confiança por \emph{bootstrap} de blocos usa um bloco base de sete ou oito dias}
sobre resíduos com autocorrelação até $0{,}90$; a sua cobertura nominal não
está demonstrada. A validação por simulação foi feita a $\rho = 0{,}90$ com
blocos de 60 dias, e não com o bloco base, e deu $357/400 = 0{,}8925$
(intervalo exato $0{,}86$ a $0{,}92$), o que não é evidência de cobertura
inferior a $0{,}90$ nem prova de cobertura de $0{,}95$. Um
contraste que exclui zero por pouco está dentro do erro de calibração do
próprio instrumento, e a sensibilidade ao comprimento de bloco reportada
nas Secções~\ref{sec:estabilidade} e~\ref{sec:calibracao} é o sintoma
visível disso. A correção de multiplicidade aplica-se só à família dos
testes de estabilidade; todos os outros resultados são por contraste, e a
sua contagem não é evidência conjunta. ``Não rejeitado'' nunca é
``equivalente''~\cite{Lakens2017}: não há margens práticas elicitadas.

\paragraph{Da simulação.} O gerador primário tem $\rho = 0{,}6$, o regime do
fuel gás, não o dos vapores; depende da relação de referência, das
máscaras, do ruído e do orçamento de réplicas escolhidos; o modelo aditivo
corre com hiperparâmetros fixos; o desvio mínimo detetável é do gerador, não
da unidade; e um número obtido em quatrocentas réplicas de um gerador que
escrevemos não é um facto sobre a refinaria, por muitas réplicas que tenha.

\paragraph{Das linhas orientadoras.} A classe linear robustificada, com
bandas que modelam a dependência e critério de aceitação, leva o modelo
estático tão longe quanto esta demonstração permitiu --- e não resolve a não-estacionariedade estrutural que o
diagnóstico documenta. O critério de aceitação é um procedimento com
critérios declarados, demonstrado com $n = 5$; o que ele decide depende do
limiar, e o que não decide é a impossibilidade da classe. As opções
de investigação seguintes --- modelos com estado de degradação, referência
congelada com monitor sequencial em produção, validação com eventos reais
--- estão no Capítulo~\ref{cha:conclusoes}; os resultados não demonstram que
modelos com estado sejam necessários nem que bastem. Para elas ficam prontos
a plataforma, o diagnóstico, o critério de aceitação e a evidência auditável
que os sustenta.
